\chapter{Image stacking}
This chapter analyze why (almost any type of) noise can be removed
pixel-wise averaging the noisy images. This denoising technique is
commonly known as \emph{image stacking}. Doing that, in theory, it is
spected to increase the \gls{SNR} by a factor of $\sqrt{N}$, where $N$
is the number of \emph{averaged}\footnote{Astronomers use the term
  ``stacked''.} images.

%Therefore, using stacking, if you were to take the same image, for
%example, 1000 times and average them together, the random fluctuations
%should cancel each other out. What remains is the mean. In other words,
%if we were able to expose (observe) the sample an infinite amount
%time, we would capture the mean of the signal (even in the presence of
%noise).

\section{Convergence}
The variance of the stacked image is the sum of the variances of each image:
\begin{equation}
  \text{Var}(\overline{\{\bm{X}^{(i)}\}_{i=0}^{N-1}}) = sum_{i=0}^{N-1}\text{Var}(\bm{X}^{(i)}).
\end{equation}
Now, if all the noisy images has the same variance
$\text{Var}(\bm{X})$, we finally obtain that
\begin{equation}
  \text{Var}(\overline{\{\bm{X}^{(i)}\}_{i=0}^{N-1}}) = N\text{Var}(\bm{X}).
\end{equation}


\subsection{The ``expected'' noisy image (Gaussian noise)}
Under the assumption that we have a high-enought number of noisy
instances $\{\tilde{\bm{X}}^{(i)}\}$ and therefore the pixel-wise
arithmetic average (see Section~\ref{sec:average}) of them can be
considered as a good estimation for the expected noisy image,
\begin{equation}
  \overline{\{\tilde{\bm{X}}^{()}\}} \approx \mathbb{E}[\{\tilde{\bm{X}}^{()}\}],
\end{equation}
the expectation for the $p$-th pixel of the noisy image is
\begin{equation}
  \mathbb{E}[\{\tilde{\bm{X}_p}\}^{()}] = \mathbb{E}[\text{max}(0, \{\bm{X}_p+\bm{N}_p)\}^{()}].
  \label{eq:E_AWG}
\end{equation}
Notice that
$\mathbb{E}[\{\tilde{\bm{X}}\}^{()}]>\mathbb{E}[\{\bm{X}\}^{()}]$ when
$\mathbb{E}[\{\bm{X}\}^{()}]=\bm{0}$. This explain the existence of
noise in \emph{dark images}\footnote{A dark image is the obtained in
  absence of radiation}. This basically implies that $\tilde{\bm{X}}$ is
an unbiased estimator for $\bm{X}$ when $\mathbb{E}[\bm{X}]>0$.

If you average $K$ independent noisy captures
$\overline{\{\tilde{\bm{X}}\}^{()}\}}$ then
\begin{equation}
  \mathbb{E}[\overline{\{\tilde{\bm{X}}\}^{()}\}}]=\bm{X},
\end{equation}
but the variance
\begin{equation}
  \text{Var}(\overline{\{\tilde{\bm{X}}\}^{()}\}}) = \frac{\sigma^2}{K},
\end{equation}
so the noise standard deviation decreases like $\sigma/\sqrt{K}$.

\section{Removing Gaussian noise}
\subsection{Consequences of the \gls{CLT}}
%\section{The mean is the clean signal}

%\section{The standard deviation is the noise}}
%"Noise" is defined as any unwanted variation that obscures the signal. In statistics, the standard deviation (σ) is the literal measurement of variation.

%    The Logic: If every pixel recorded the exact same value (μ) every time, there would be zero noise. But because of the random nature of particles, the values "wiggle" around the mean.

\subsection{Applicability}
%{{{

Stacking relies on the (Lindeberg-Levy version of the) \gls{CLT},
which is applicable only if:
\begin{enumerate}
\item The mean of the noise is finite: $\mu<\infty$.
\item The variance of the noise is finite: $\sigma^2\infty$.
\end{enumerate}
Fortunatelly, most of the natural noise sources satisfy these
conditions.\footnote{However, notice that there are statistical
  distributions, such as the Cauchy distribution that has infinite
  variance and an undefined mean. Thus, the addition of a infinite
  number of sequences of random number that follow (each of them) a
  Cauchy distribution is the Cauchy distribution.} Futhermore, there
exists a more relaxed ``version'' of the \gls{CLT} known as the
Lindeberg-Feller \gls{CLT} which states that the \gls{CLT} it is
satisfied even if the random variables are not \gls{iid}, provided
that no single variable in the sequence dominates the others. This is
why the sum of many different types of noise (such as thermal noise +
shot noise + interference) still ends up looking Gaussian (see
Section~\ref{sec:CLT}).

Therefore, under a high enough number of noisy instancces, the
\gls{CLT} allows us to take into consideration only the \gls{AWG}
noise case (see Section~\ref{sec:AWG}) when we perform denoising
throught computing the mean of those noisy instances. Let's
\begin{equation}
  \hat{\bm{x}}^{(i)} = \bm{x} + \bm{n}^{(i)},
  \label{eq:noisy_instances}
\end{equation}
represents the generation of the $i$-th noisy instance of the signal
$\bm{x}$. Assuming that
\begin{equation}
  \mathbb{E}(\bm{n})=\bm{0},
\end{equation}
where $\bm{n}=\{\bm{n}^{(i)}\}_{i=0}^{N-1}$ is the set of
random signals, and taking expectations in
Eq.~\ref{eq:noisy_instances}, we get
\begin{equation}
  \mathbb{E}(\hat{\bm{x}}) = \mathbb{E}(\bm{x}) + \mathbb{E}(\bm{n}) = \bm{x},
  \label{eq:MNI}
\end{equation}
where $\hat{\bm{x}}=\{\hat{\bm{s}}^{(i)}\}_{i=0}^{N-1}$.

Notice, however, that Eq.~\ref{eq:MNI} is only true when
$N\rightarrow\infty$. In practice $N$ is finite, and we only get an
approximation $\tilde{\bm{s}}^{[N]}$, where $N$ is the number of
averaged signals.

%}}}


\subsection{\gls{SNR}}
The \gls{SNR} is defined as the ratio between the power (see
Eq.~\ref{eq:power_as_expectation}) of the clean signal and the power of
the noise:
\begin{equation}
  \text{SNR} = \frac{P(\bm{x})}{P(\bm{n})} = \frac{\mathbb{E}[\bm{x}^2]}{\mathbb{E}[\bm{n}^2]} = \frac{\mathbb{E}[\bm{x}]^2+\text{Var}(\bm{x})}{\mathbb{E}[\bm{n}]^2+\text{Var}(\bm{n})}
\end{equation}
Considering now that for the \gls{AWG} model, $\mathbb{E}[\bm{n}] = 0$,
\begin{equation}
  \text{SNR} = \frac{\mathbb{E}[\bm{x}]^2+\text{Var}(\bm{x})}{\sigma^2}
\end{equation}

\subsection{SNR estimation}
Unfortunately, when only the noisy signal is available, the \gls{SNR}
only can be estimated. If the clean signal is concentrated in the
frequency domain, a procedure to do this is:
\begin{enumerate}
\item Compute $\mathcal{F}(\hat{\bm{x}})$.
\item Identify the ``noise floor'' -- the flat part of the spectrum
  where no signal peaks exist.
\item Estimate $P(\bm{n})$ as the average level of that flat
  floor. Notice that in the case of Gaussian noise, you have estimated
  $\sigma^2$.
\item Estimate $P(\bm{x})$ by subtracting the noise floor from the
  total power of the peaks.
\end{enumerate}

%\subsection{\acrshort{SNR} of the \acrshort{AWG} noise model}

%The \gls{SNR} is commonly defined as the ratio between the power (see
%Eq.~\ref{eq:power_as_expectation} and Eq.~\ref{eq:variance}) of the
%clean signal and the power of the noise:
%\begin{equation}
%  \text{SNR} = \frac{P(\mathbf{s})}{P(\mathbf{n})} = \frac{\overline{\mathbb{E}}(\mathbf{s}^2)}{\overline{\mathbb{E}}(\mathbf{n}^2)} = \frac{\overline{\mathbb{E}}(\mathbf{s})^2+\overline{\mathbb{V}}(\mathbf{s})}{\overline{\mathbb{E}}(\mathbf{n})^2+\overline{\mathbb{V}}(\mathbf{n})},
%  \label{eq:SNR}
%\end{equation}
%where $\mathbf{s}$ is the clean signal and $\hat{\mathbf{s}}$ is its
%noisy version, and
%\begin{equation}
%  \mathbf{n} = \hat{\mathbf{s}} - \mathbf{s}.
%\end{equation}

%Considering that $\overline{\mathbb{E}}(\mathbf{n}) = 0$,
%\begin{equation}
%  \text{SNR} = \frac{\overline{\mathbb{E}}(\mathbf{s})^2 + \overline{\mathbb{V}}(\mathbf{s})}{\sigma^2},
%\end{equation}

\begin{comment}
that
\begin{equation}
  \mathbb{E}(\mathbf{s}) = \mathbf{s},
  \label{eq:E_constant}
\end{equation}
that $\hat{\mathbf{s}}$ is an unbiased estimator of $\mathbf{s}$, i.e.,
\begin{equation}
  \lim_{N\rightarrow\infty}\frac{1}{N}\sum_{n=0}^{N-1}\hat{\mathbf{s}}^{(n)} = \mathbf{s},
\end{equation}
where $\hat{\mathbf{s}}^{(n)}$ is the $n$-th noisy instance of
$\mathbf{s}$, and using Eq.~\ref{eq:SNR_distribution}, we come to
\begin{equation}
  \text{SNR} = \frac{\textbf{s}}{\sigma} = \frac{\hat{\textbf{s}}}{\sigma},
  \label{eq:SNR_AWG}
\end{equation}
and therefore to
\begin{equation}
  \text{SNR}_i = \frac{\hat{\textbf{s}}_i}{\sigma}.
  \label{eq:SNR_AWG_i}
\end{equation}
\end{comment}

\subsection{Average convergence of random numbers}
The \gls{AWG} noise model is used in most denoising algorithms because
the component-wise arithmetic average\footnote{In general, the
  addition.} of sequences of random numbers always converges to a
normal distribution (see Section \ref{sec:CLT}). In fact, a common way
of generating a sequence of close-normal random numbers is:
\begin{enumerate}
\item Generate 12 sequences of uniform random numbers between
  $[0, 1[$. The $\mathcal{U}(0,1)$ has $\mu=0.5$ and $\sigma^2=1/12$.
\item Add them, component-wise, obtaining a new sequence that has
  $\mu=6$ and $\sigma^2=1$.
\item Subtract $6$, to have $\mu=0$.
\end{enumerate}
Notice that the higher the number of noise realizations added, the
higher the accuracy of this algorithm.

\section{Averaging multiple acquisitions}
If you take $K$ independent Poisson measurements
$\{\tilde{\bm{X}}^{(k)}\}$, where
\begin{equation}
  \tilde{\bm{X}}_p^{(k)}\sim\mathcal{P}(\bm{X}_p)\forall p
\end{equation}
and average them pixelwise, the expectation for $p$ is
\begin{equation}
  \mathbb{E}[\{\tilde{\bm{X}}_p^{(k)}\}] = \bm{X}_p
\end{equation}
with variance
\begin{equation}
  \text{Var}(\{\tilde{\bm{X}}_p^{(k)}\}) = \frac{\bm{X}_p}{K}
\end{equation}
and the A-\gls{SNR} improves as $\sqrt{K\bm{X}_p}$.

\section{Experiments (MPG)}

This section is basically an experimental verification of the
efficiency of the \gls{CLT} (see
Section~\ref{sec:CLT}). Figure~\ref{fig:MNI_0MMPG} shows an example of
how \gls{MNI} increases the quality using images with \gls{MPG} noise.
  
\begin{figure}
  \centering
  \resizebox{1.0\textwidth}{!}{
    \renewcommand{\arraystretch}{0.0} % Adjust row spacing in the table
    \setlength{\tabcolsep}{0ex}      % Adjust column spacing in the table    
    \begin{tabular}{cc}
      \href{https://nbviewer.org/github/vicente-gonzalez-ruiz/denoising/blob/main/notebooks/barb.ipynb\#barb}{\includegraphics{barb.pdf}} & \href{https://nbviewer.org/github/vicente-gonzalez-ruiz/denoising/blob/main/notebooks/barb_0MMPG.ipynb\#barb_0MMPG}{\includegraphics{barb_0MMPG.pdf}} \\
      \href{https://nbviewer.org/github/vicente-gonzalez-ruiz/denoising/blob/main/notebooks/barb_averaging.ipynb\#barb_averaging}{\includegraphics{barb_averaging.pdf}} & \href{https://nbviewer.org/github/vicente-gonzalez-ruiz/denoising/blob/main/notebooks/barb_averaging_PSNR.ipynb\#barb_averaging_PSNR}{\includegraphics{barb_averaging_PSNR.pdf}}
    \end{tabular}
  }
  \caption{Increase of the SNR as a function of the number of noisy
    instances used in MNI. The clean image is shown on the top-left,
    and a noisy instance on the top-right. On the bottom-left, the
    denoised NMI image, and on the bottom-right the progression of the
    SNR as a function of $I$ and the level of
    noise.\label{fig:MNI_0MMPG}}
\end{figure}

\begin{comment}

Fig.~\ref{fig:0MAUN} shows an
example of how zero-mean uniform noise is cancelled by averaging.
  
  \begin{figure}
    \centering
    \resizebox{1.0\textwidth}{!}{
      \renewcommand{\arraystretch}{0.0} % Adjust row spacing in the table
      \setlength{\tabcolsep}{0ex}      % Adjust column spacing in the table    
      \begin{tabular}{cc}
        \href{https://nbviewer.org/github/vicente-gonzalez-ruiz/denoising/blob/main/figs/averaging_denoising.ipynb\#Display-Barb}{\includegraphics{barb}} & \href{https://nbviewer.org/github/vicente-gonzalez-ruiz/denoising/blob/main/figs/averaging_denoising.ipynb\#0MAUN_barb}{\includegraphics{0MAUN_barb}} \\
        \href{https://nbviewer.org/github/vicente-gonzalez-ruiz/denoising/blob/main/figs/averaging_denoising.ipynb\#denoised_0MAUN_barb}{\includegraphics{denoised_0MAUN_barb}} & \href{https://nbviewer.org/github/vicente-gonzalez-ruiz/denoising/blob/main/figs/averaging_denoising.ipynb\#PSNR_0MAUN_barb}{\includegraphics{PSNR_0MAUN_barb}}
      \end{tabular}
    }
    \caption{Effect of zero-mean additive uniform noise in an image
      and how averaging can be used to reduced by averaging. The clean
      image of Barb is shown on the top left, and a noisy version on
      the top right. On to bottom, the left image shows a denoised
      version after averaging and the right graph shows the performance
      of the averaging process for different levels of
      noise.\label{fig:0MAUN}}
  \end{figure}

 Fig.~\ref{fig:0MAGN} shows an
  example of how zero-mean Gaussian noise is cancelled by averaging.

  \begin{figure}
    \centering
    \resizebox{1.0\textwidth}{!}{
      \renewcommand{\arraystretch}{0.0} % Adjust row spacing in the table
      \setlength{\tabcolsep}{0ex}      % Adjust column spacing in the table    
      \begin{tabular}{cc}
        \href{https://nbviewer.org/github/vicente-gonzalez-ruiz/denoising/blob/main/figs/averaging_denoising.ipynb\#barb}{\includegraphics{barb}} & \href{https://nbviewer.org/github/vicente-gonzalez-ruiz/denoising/blob/main/figs/averaging_denoising.ipynb\#0MAGN_barb}{\includegraphics{0MAGN_barb}} \\
        \href{https://nbviewer.org/github/vicente-gonzalez-ruiz/denoising/blob/main/figs/averaging_denoising.ipynb\#denoised_0MAGN_barb}{\includegraphics{denoised_0MAGN_barb}} & \href{https://nbviewer.org/github/vicente-gonzalez-ruiz/denoising/blob/main/figs/averaging_denoising.ipynb\#PSNR_0MAGN_barb}{\includegraphics{PSNR_0MAGN_barb}}
      \end{tabular}
    }
    \caption{Effect of zero-mean additive Gaussian noise in an image
      and how averaging can be used to reduced by averaging. The clean
      image of Barb is shown on the top left, and a noisy version on
      the top right. On to bottom, the left image shows a denoised
      version after averaging and the right graph shows the performance
      of the averaging process for different levels of
      noise.\label{fig:0MAGN}}
  \end{figure}

  Fig.~\ref{fig:0MMGN} shows an example of
  how zero-mean Gaussian noise is cancelled by averaging.

  \begin{figure}
    \centering
    \resizebox{1.0\textwidth}{!}{
      \renewcommand{\arraystretch}{0.0} % Adjust row spacing in the table
      \setlength{\tabcolsep}{0ex}      % Adjust column spacing in the table    
      \begin{tabular}{cc}
        \href{https://nbviewer.org/github/vicente-gonzalez-ruiz/denoising/blob/main/figs/averaging_denoising.ipynb\#barb}{\includegraphics{barb}} & \href{https://nbviewer.org/github/vicente-gonzalez-ruiz/denoising/blob/main/figs/averaging_denoising.ipynb\#0MMGN_barb}{\includegraphics{0MMGN_barb}} \\
        \href{https://nbviewer.org/github/vicente-gonzalez-ruiz/denoising/blob/main/figs/averaging_denoising.ipynb\#denoised_0MMGN_barb}{\includegraphics{denoised_0MMGN_barb}} & \href{https://nbviewer.org/github/vicente-gonzalez-ruiz/denoising/blob/main/figs/averaging_denoising.ipynb\#PSNR_0MMGN_barb}{\includegraphics{PSNR_0MMGN_barb}}
      \end{tabular}
    }
    \caption{Effect of zero-mean multiplicative Gaussian noise in an
      image and how averaging can be used to reduced by averaging. The
      clean image of Barb is shown on the top left, and a noisy
      version on the top right. On to bottom, the left image shows a
      denoised version after averaging and the right graph shows the
      performance of the averaging process for different levels of
      noise.\label{fig:0MMGN}}
  \end{figure}

 Fig.~\ref{fig:Rayleigh}
  shows an example of how Rayleigh noise is cancelled by averaging.

  \begin{figure}
    \centering
    \resizebox{1.0\textwidth}{!}{
      \renewcommand{\arraystretch}{0.0} % Adjust row spacing in the table
      \setlength{\tabcolsep}{0ex}      % Adjust column spacing in the table    
      \begin{tabular}{cc}
        \href{https://nbviewer.org/github/vicente-gonzalez-ruiz/denoising/blob/main/figs/averaging_denoising.ipynb\#barb}{\includegraphics{barb}} & \href{https://nbviewer.org/github/vicente-gonzalez-ruiz/denoising/blob/main/figs/averaging_denoising.ipynb\#Rayleigh_barb}{\includegraphics{Rayleigh_barb}} \\
        \href{https://nbviewer.org/github/vicente-gonzalez-ruiz/denoising/blob/main/figs/averaging_denoising.ipynb\#denoised_Rayleigh_barb}{\includegraphics{denoised_Rayleigh_barb}} & \href{https://nbviewer.org/github/vicente-gonzalez-ruiz/denoising/blob/main/figs/averaging_denoising.ipynb\#PSNR_Rayleigh_barb}{\includegraphics{PSNR_Rayleigh_barb}}
      \end{tabular}
    }
    \caption{Effect of Rayleigh noise in an image and how averaging can
      be used to reduced by averaging. The clean image of Barb is shown
      on the top left, and a noisy version on the top right. On to
      bottom, the left image shows a denoised version after averaging
      and the right graph shows the performance of the averaging
      process for different levels of noise.\label{fig:Rayleigh}}
  \end{figure}

  Fig.~\ref{fig:Poisson} shows an example of how Poisson noise is
  cancelled by averaging.

  \begin{figure}
    \centering
    \resizebox{1.0\textwidth}{!}{
      \renewcommand{\arraystretch}{0.0} % Adjust row spacing in the table
      \setlength{\tabcolsep}{0ex}      % Adjust column spacing in the table    
      \begin{tabular}{cc}
        \href{https://nbviewer.org/github/vicente-gonzalez-ruiz/denoising/blob/main/figs/averaging_denoising.ipynb\#barb}{\includegraphics{barb}} & \href{https://nbviewer.org/github/vicente-gonzalez-ruiz/denoising/blob/main/figs/averaging_denoising.ipynb\#Poisson_barb}{\includegraphics{Poisson_barb}} \\
        \href{https://nbviewer.org/github/vicente-gonzalez-ruiz/denoising/blob/main/figs/averaging_denoising.ipynb\#denoised_Poisson_barb}{\includegraphics{denoised_Poisson_barb}} & \href{https://nbviewer.org/github/vicente-gonzalez-ruiz/denoising/blob/main/figs/averaging_denoising.ipynb\#PSNR_Poisson_barb}{\includegraphics{PSNR_Poisson_barb}}
      \end{tabular}
    }
    \caption{Effect of Poisson noise in an image and how averaging can
      be used to reduced by averaging. The clean image is shown
      on the top left, and a noisy version on the top right. On the
      bottom, the left image shows a denoised version after averaging
      and the right graph shows the performance of the averaging
      process for different levels of noise.\label{fig:Poisson}}
  \end{figure}

\end{comment}

